% ==============================================================================
% T0/FFGFT Theory: Document 221
% Title: Redshift falsification criteria:
%        Observational signatures beyond H_0 reproduction
% Author: Johann Pascher
% Date: May 2026
% References: Documents 026 (Geometric Cosmology), 064 (Hubble constant),
%             165 (H_0 ratio), 220 (Casimir falsification)
% Occasion: P. Austin (Information Physics Institute), falsification request
% ==============================================================================

\section*{Preamble}

This document formulates experimentally testable falsification criteria
for the FFGFT prediction on cosmological redshift, in direct response
to P.~Austin's methodological question:
\emph{\enquote{What observational signature distinguishes geometric
path-stretch from metric expansion, beyond reproducing an $H_0$-like
number?}}

FFGFT is \emph{not} a standard \enquote{tired light} model. The four
classical objections to tired light (frequency-dependent scattering,
absent time dilation, Tolman test, CMB spectrum) are specifically
addressed in the FFGFT prediction. This document shows how -- and gives
the concrete falsification criterion in each case.

% ==============================================================================
\section{\texorpdfstring{What FFGFT predicts: $\xi$-vacuum energy loss}{What FFGFT predicts: xi-vacuum energy loss}}
% ==============================================================================

From Doc.~064, the FFGFT mechanism is:
\begin{equation}
  \frac{dE}{dx} = -\xi_0 \cdot E,
  \qquad
  E(x) = E_0 \cdot e^{-\xi_0 x},
  \label{eq:energy-loss}
\end{equation}
where $x$ is the travelled distance in suitable units. For
cosmological distances ($\xi_0 x \ll 1$):
\begin{equation}
  z \approx \xi_0 \cdot x.
\end{equation}

From Doc.~165 the numerical value follows:
\begin{equation}
  H_0^{\text{T0}}
  = \frac{E_H}{\hbar},
  \qquad
  E_H = E_0 \cdot \xi_0^{41/4},
  \qquad
  H_0^{\text{T0}} \approx 66.2\,\text{km/s/Mpc}.
\end{equation}

This value clearly favours the \emph{Planck value} ($H_0 \approx 67$),
not the local SH0ES value ($H_0 \approx 73$). It is not a fit, but a
\textbf{completely parameter-free} prediction from $\xi_0 = 4/30000$
and the electron mass.

% ==============================================================================
\section{\texorpdfstring{Four observational signatures beyond $H_0$}{Four observational signatures beyond H_0}}
% ==============================================================================

\subsection{Signature 1: frequency-independence of redshift}

\textbf{Standard tired light fails here:} scattering models (Compton,
plasma) produce \emph{frequency-dependent} redshift, since the cross
section depends on photon energy. Observationally this is \emph{not}
the case: quasar spectra at $z = 7$ show identical line shifts for UV
and radio photons.

\textbf{FFGFT prediction:} the $\xi_0$ energy loss is a
\emph{geometric} property of spacetime itself, not a scattering
interaction. From (\ref{eq:energy-loss}) the relative energy loss
$\Delta E/E = -\xi_0 \cdot dx$ is identical for \emph{all} photon
energies -- the model is intrinsically frequency-independent.

\textbf{Falsification:} a systematic frequency dependence
$|\Delta z(\nu_1) - \Delta z(\nu_2)| / z > 10^{-4}$ over several
orders of magnitude in $\nu$ would refute FFGFT. Current observations
(Sloan, JWST) show $< 10^{-5}$ frequency dependence, consistent with
FFGFT.

\subsection{Signature 2: Hubble tension as internal prediction}

\textbf{Standard cosmology has a problem:} local distance-ladder
measurements (SH0ES, Cepheids + type-Ia supernovae) give
$H_0 = 73.0 \pm 1.0$ km/s/Mpc; CMB fitting (Planck) gives
$H_0 = 67.4 \pm 0.5$ km/s/Mpc. The discrepancy of $\sim 9\,\%$ is
highly significant ($> 5\sigma$) and unresolved.

\textbf{FFGFT prediction:} the \emph{true} $H_0$ is
$H_0^{\text{T0}} \approx 66.2$ km/s/Mpc, close to the Planck value.
The discrepancy with SH0ES is not new physics, but a methodological
issue: local distance-ladder methods rely on expansion assumptions
(Hubble flow, peculiar-velocity corrections) that are systematically
biased in FFGFT. Specifically:
\begin{itemize}[leftmargin=2em,itemsep=2pt,topsep=2pt]
  \item The SH0ES method calibrates Cepheids in \emph{local}
        galaxies (distances up to $\sim 40$ Mpc), where the
        $\xi_0$ energy loss gives $\xi_0 x \approx 5\!\cdot\!10^{-3}$
        -- small, but not zero.
  \item The cosmological CMB method measures over $z \sim 1100$ and
        is robust against local effects.
  \item In FFGFT the apparent excess of $\sim 9\,\%$ in the local
        $H_0$ determination arises as a consequence of the static
        geometry: the Hubble law $z = H_0 d/c$ holds with a different
        constant for local vs.\ cosmological scales, because the
        distance measurement itself is $\xi_0$-affected.
\end{itemize}

\textbf{Falsification:} if independent, expansion-free distance
measurements (e.g.\ purely geometric methods via maser galaxies or
gravitational-wave standard sirens) systematically yield
$H_0 \neq 66.2 \pm 2$ km/s/Mpc, FFGFT is refuted. \emph{Currently:}
NGC~4258 masers give $H_0 = 72.0 \pm 3$ km/s/Mpc (SH0ES-consistent,
FFGFT-tense), GW standard sirens give $H_0 = 70 \pm 7$ km/s/Mpc (too
uncertain to discriminate). Future LIGO-Virgo-KAGRA data and
JWST-NGC~4258 improvements are the critical tests.

\subsection{Signature 3: supernova time dilation}

\textbf{This is the hardest classical objection to tired light:}
high-redshift type-Ia supernovae show light curves \emph{broadened}
by the factor $(1+z)$ (Goldhaber et al.\ 2001; Riess et al.\ 1997).
In standard expansion models this follows from metric time dilation.
In standard tired-light models it is unexplained.

\textbf{FFGFT prediction:} in FFGFT the apparent time dilation is
\emph{also} present, but from a geometric rather than metric source.
The argument:

\begin{itemize}[leftmargin=2em,itemsep=2pt,topsep=2pt]
  \item Photons from the start of a supernova light curve and photons
        from the end traverse the $\xi_0$ path equally long
        \emph{in distance}, but the observed time interval between
        them is stretched by the factor $(1+z) = 1/(1 - \xi_0 x)$,
        because the relative energy loss also stretches the
        \emph{frequency domain}.
  \item More precisely: if a source emits photons with period $T_0$,
        we observe after path-stretching the period
        $T_{\text{obs}} = T_0 \cdot (1+z)$, indistinguishable from
        metric time dilation at this order.
\end{itemize}

\textbf{Falsification:} a supernova-light-curve analysis showing
broadening that does \emph{not} follow $(1+z)$ scaling, but a
different law (e.g.\ constant broadening independent of $z$, as pure
scattering would predict), would refute both standard expansion and
FFGFT. Currently all data show $(1+z)$ scaling, consistent with both.

\textbf{Discriminating test:} at very high redshifts ($z > 2$),
$(1+z)$ is \emph{no longer} identical with $1/(1 - \xi_0 x)$, since
the linear approximation in (\ref{eq:energy-loss}) breaks. FFGFT says:
\begin{equation}
  T_{\text{obs}}/T_0 = e^{\xi_0 x} = 1 + z + \tfrac{1}{2}z^2 + \ldots
\end{equation}
while standard FRW gives:
\begin{equation}
  T_{\text{obs}}/T_0 = 1 + z.
\end{equation}
At $z = 2$ the difference is $z^2/2 = 2$, hence a factor of 1.5 in
light-curve width. \textbf{This is measurable and unambiguously
falsifies FFGFT or standard cosmology.} JWST observations of $z > 2$
supernovae are the key test.

\subsection{Signature 4: Tolman surface-brightness test}

\textbf{Standard tired light fails here:} a static universe gives
surface brightness $\propto (1+z)^{-1}$, an expanding one
$\propto (1+z)^{-4}$. Observations (Lubin \& Sandage 2001 for $z < 1$;
Pahre et al.\ 1996; recent work to $z = 5$) show $(1+z)^{-4}$.

\textbf{FFGFT prediction:} in FFGFT the static geometry is not simply
\enquote{no effect}, but contains the $\xi_0$ energy loss with
frequency-independence \emph{and} a geometric bunching component
acting analogously to the $(1+z)^{-2}$ bundle stretching in expanding
models:
\begin{itemize}[leftmargin=2em,itemsep=2pt,topsep=2pt]
  \item Energy loss per photon: factor $(1+z)^{-1}$ (from $E \propto e^{-\xi_0 x}$)
  \item Photon arrival rate: factor $(1+z)^{-1}$ (from time dilation, signature 3)
  \item Geometric path bunching: factor $(1+z)^{-2}$ (from
        $\xi_0$-vacuum granularity, Doc.~091/220)
\end{itemize}
Product: $(1+z)^{-4}$, identical to the standard prediction.

\textbf{This is FFGFT's distinguishing statement vis-à-vis tired
light:} the geometric path-bunching component is FFGFT-specific
(follows from the same $\xi_0$-vacuum structure that produces the
Casimir effect) and is not present in simple tired-light models.

\textbf{Falsification:} if the Tolman test at very high $z$
($z > 5$, only accessible with JWST) systematically deviates from
$(1+z)^{-4}$ and shows a \emph{simpler} scaling exponent (e.g.\
$(1+z)^{-3}$ or $(1+z)^{-1}$), FFGFT is refuted -- this would mean
the geometric path bunching does not exist in the predicted form.

% ==============================================================================
\section{Specific falsification criteria}
% ==============================================================================

FFGFT is \textbf{refuted} by any of the following observations:

\begin{enumerate}[leftmargin=2em,itemsep=4pt,topsep=2pt]
  \item \textbf{Frequency dependence:}
        a measurable ($> 10^{-4}$) frequency dependence of the
        cosmological redshift over several orders of magnitude in $\nu$.
        \emph{Current status:} no indication, FFGFT consistent.

  \item \textbf{$H_0$ value:}
        a future expansion-independent $H_0$ measurement
        (GW standard sirens with $\sigma < 1$ km/s/Mpc, maser methods)
        yielding $H_0 \neq 66.2 \pm 2$ km/s/Mpc.
        \emph{Current status:} undecided, critical test in next 5 years.

  \item \textbf{High-$z$ supernova time dilation:}
        light-curve broadening at $z > 2$ following $(1+z)$ rather
        than $e^{\xi_0 x} = 1 + z + z^2/2 + \ldots$
        \emph{Current status:} JWST data just becoming available;
        critical test for the next few years.

  \item \textbf{Tolman scaling at $z > 5$:}
        surface-brightness scaling that significantly deviates from
        $(1+z)^{-4}$ (in particular $(1+z)^{-3}$ or $(1+z)^{-1}$).
        \emph{Current status:} JWST data to $z = 12$ available;
        galaxy evolution makes these tests model-dependent.

  \item \textbf{CMB spectrum:}
        a deviation of the CMB spectrum from a perfect Planck
        blackbody at $T = 2.725$ K with $|\Delta T/T| > 10^{-4}$
        unexplained by known systematics.
        \emph{Current status:} COBE/FIRAS finds
        $|\Delta T/T| < 10^{-5}$, FFGFT consistent (see Doc.~091
        for the CMB-Casimir connection).
\end{enumerate}

% ==============================================================================
\section{What FFGFT does NOT predict}
% ==============================================================================

For methodological clarity:

\begin{itemize}[leftmargin=2em,itemsep=2pt,topsep=2pt]
  \item FFGFT does \emph{not} predict any violation of the photon
        speed $c$. Photons propagate on null geodesics in the
        $\xi_0$-vacuum geometry.
  \item FFGFT does \emph{not} predict any redshift anisotropy not
        explainable by known local motion. The $\xi_0$-vacuum effect
        is spatially isotropic.
  \item FFGFT does \emph{not} predict any time variation of $\xi_0$
        -- it is a fundamental dimensionless geometric parameter,
        not an evolving field. Variation tests of natural constants
        (quasar spectra at high $z$) should therefore find
        \emph{no} variation.
        \emph{Current status:} no variation found, FFGFT consistent.
\end{itemize}

% ==============================================================================
\section{Connection to Doc.~220 (Casimir falsification)}
% ==============================================================================

The $\xi_0$-vacuum mechanism producing the cosmological redshift is
\emph{the same} mechanism producing the Casimir effect. Both effects
are linked by the bridge formula
\begin{equation}
  \rho_\text{CMB} = \frac{\xi_0 \hbar c}{L_\xi^4},
  \qquad
  L_\xi \approx 100\,\mu\text{m}
\end{equation}
(Doc.~091, 166). This implies \emph{consistent} predictions:

\begin{itemize}[leftmargin=2em,itemsep=2pt,topsep=2pt]
  \item Casimir falsification (Doc.~220) tests
        $\xi_0$-vacuum on the $\sim 100\,\mu$m scale.
  \item Redshift falsification (this document) tests the
        \emph{same} $\xi_0$ mechanism on cosmological scales.
  \item Falsification at one scale implies falsification at the
        other, since both follow from the same vacuum geometry.
\end{itemize}

\textbf{Hierarchical falsification.} If the Casimir prediction
(Doc.~220) is refuted at $d \to L_\xi$, the redshift prediction is
also affected, since $L_\xi$ enters via $\rho_\text{CMB}$ in both
predictions. Conversely, if the redshift prediction (e.g.\
$H_0 \neq 66.2$) is refuted, the $\xi_0$-vacuum geometry itself is
called into question, retroactively destroying the Casimir prediction.
The two tests are therefore \emph{not} independent, but probe
different scale-manifestations of the same phenomenon.

% ==============================================================================
\section{Summary}
% ==============================================================================

FFGFT makes the following quantitative, falsifiable predictions in
the cosmological domain:

\begin{enumerate}[leftmargin=2em,itemsep=2pt,topsep=2pt]
  \item $H_0 = (66.2 \pm 1.3)$ km/s/Mpc
        (parameter-free from $\xi_0 = 4/30000$ and $m_e$).
  \item Frequency-independence of redshift to $< 10^{-4}$.
  \item Supernova time dilation $T_{\text{obs}}/T_0 = e^{\xi_0 x}$,
        distinguishable from standard $(1+z)$ at $z > 2$.
  \item Tolman scaling $(1+z)^{-4}$ from three combined effects
        (energy loss, time dilation, geometric bunching), not from
        expansion.
  \item No variation of $\xi_0$ over cosmic time.
\end{enumerate}

\textbf{Methodological note.} These predictions are FFGFT-internal.
They follow from the $\xi_0$-vacuum energy loss and the $\xi_0$-vacuum
geometry, not from topological partitions or recursive survival
structures. The prediction $H_0 \approx 66.2$ km/s/Mpc is consistent
with the Planck value and in conflict with the SH0ES value -- the
Hubble tension is systematically explained by the methodological
contamination of expansion-based distance ladders.

\textbf{Status.} Three of the five tests (frequency-independence, CMB
spectrum, constancy of $\xi_0$) are currently consistent with FFGFT.
Two tests (supernova time dilation at $z > 2$, Tolman at $z > 5$)
become critical in the coming years with JWST data. An
expansion-free $H_0$ measurement with $< 1\,\%$ accuracy (e.g.\
through GW standard sirens) is the sharpest discriminator.
